\documentclass[11pt, a4paper]{article}

\usepackage[utf8]{inputenc}
\usepackage[T1]{fontenc}
\usepackage{lmodern}
\usepackage[margin=2.5cm]{geometry}
\usepackage{amsmath, amssymb, amsthm}
\usepackage{booktabs}
\usepackage{enumitem}
\usepackage{hyperref}
\usepackage{xcolor}
\usepackage{microtype}
\usepackage{graphicx}
\usepackage{tikz}
\usetikzlibrary{positioning,calc,fit,arrows.meta,backgrounds}

\hypersetup{
    colorlinks=true, linkcolor=black,
    citecolor=blue!60!black, urlcolor=blue!60!black
}

\newcommand{\M}{\mathbf{M}}
\newcommand{\R}{\mathbb{R}}
\newcommand{\bigO}{\mathcal{O}}

\theoremstyle{plain}
\newtheorem{observation}{Observation}
\theoremstyle{remark}
\newtheorem*{remark}{Remark}

\title{%
    \textbf{The Orbital Response Network}\\[0.3em]
    \large From Shared Coupling Geometry to Composable Dynamics%
}
\author{Sirsh}
\date{March 2026}

\begin{document}
\maketitle

\begin{abstract}
Transformer attention coupling, the rule determining which tokens attend to which, is redundant across layers.
A spectral analysis of five pretrained model families shows that the eigenvalue distribution of $W_Q^\top W_K$ is nearly identical at every depth (correlation $0.97$) while the raw matrices are unrelated (cosine $< 0.01$).
The coupling \emph{shape} is invariant; only the \emph{orientation} varies, rotated by the per-layer response heads' updates to the residual stream.

We exploit this by sharing a single coupling matrix $\M = AB^\top$ across all layers, keeping response heads per-layer.
We call this the \textbf{Orbital Response Network (ORN)}.

At 91M parameters (V1), the ORN matches GPT-2 Small on OpenWebText (val loss 3.35) using just 524K coupling parameters (0.6\% of the model).
$\M$'s effective rank crystallises to 19--20 by 250M tokens, 10\% of training.
Freezing this $\M$ and transferring to FineWeb-Edu, the model tracks within 2\% of full training (val ${\sim}3.40$), confirming that the 10K numbers encoding relational grammar are dataset-invariant.
Multimodal experiments show sub-linear rank scaling: 4 modalities require only rank 4.

Mathematically, $\M$ is a finite-dimensional Koopman operator: its eigenvalues encode timescales, its eigenvectors encode independent dynamical modes.
We extend the ORN to the \textbf{Composable Orbital Response Network (CORN)}, where time-scaled semigroup dynamics ($t \cdot A$, $t \cdot B$) enable both prediction and compositional reasoning with no new parameters.
\end{abstract}

% ======================================================================
\section{Introduction: The Residual Stream Is the Scaling Mechanism}

The transformer~\cite{vaswani2017attention} processes sequences through two interleaved operations: attention (which determines how tokens exchange information) and feed-forward networks (which transform the exchanged information).
Both contribute to the residual stream, a persistent vector per token that accumulates updates from every layer.

We argue that these two components play fundamentally different roles:
\begin{description}[nosep, leftmargin=1.5em]
    \item[The residual stream] is the scaling mechanism.
    It is the shared medium through which all layers communicate, the surface on which representations are accumulated, and the object that persists from embedding to prediction.
    Without the residual connection, deep transformers cannot train.
    With it, information from layer~1 is directly available at layer~$L$, enabling the depth that makes transformers powerful.

    \item[The attention blocks] are the coupling mechanism.
    They determine \emph{which} tokens should exchange information at each layer, the relational structure.
    They unlock the residual stream's scaling capacity by directing information flow, but they are not themselves the source of that capacity.
\end{description}

This distinction suggests an architectural question: if the residual stream is the persistent medium, can the coupling rule that governs attention be part of that medium, shared across layers like the stream itself, rather than stored independently at each layer?

\subsection{The Cost of Per-Layer Coupling}

In a standard transformer with $L$ layers, each layer stores its own query and key projection matrices $W_Q^{(\ell)}, W_K^{(\ell)} \in \R^{d \times d}$.
The total coupling cost is $2Ld^2$ parameters, linear in depth.
This is the price of allowing each layer to define its own relational structure independently.

At the scales where transformers excel ($L = 24$--$96$, $d = 1024$--$4096$), this coupling budget is substantial: for GPT-2 XL ($L{=}48$, $d{=}1600$), the coupling parameters total ${\sim}245$M out of 1.5B, or 16\% of the model.
Each layer independently learns ``which tokens should attend to which,'' with no mechanism for sharing this knowledge across depths.

\subsection{The Spectral Invariance Observation}

Is this independence necessary?
We measured the effective coupling operator $M_\ell = W_Q^{(\ell)\top} W_K^{(\ell)}$ at every layer of five pretrained transformer families, then compared their singular value decompositions:

\begin{table}[h]
\centering\small
\begin{tabular}{@{}lrccc@{}}
\toprule
\textbf{Model} & \textbf{Params} & \textbf{Layers} & \textbf{Spectrum corr.} & \textbf{Raw cosine} \\
\midrule
GPT-2 Small & 124M & 12 & 0.982 & 0.000 \\
GPT-2 Medium & 345M & 24 & 0.962 & 0.003 \\
GPT-2 Large & 774M & 36 & 0.987 & 0.006 \\
SmolLM2 & 135M & 30 & 0.958 & 0.005 \\
Qwen2.5 & 500M & 24 & 0.940 & 0.005 \\
\midrule
\textbf{Average} & & & \textbf{0.966} & \textbf{0.004} \\
\bottomrule
\end{tabular}
\caption{Spectral invariance of $W_Q^\top W_K$ across layers. Every model shows $>0.93$ spectrum correlation with $<0.01$ raw cosine similarity. The coupling shape is invariant; the orientation varies.}
\label{tab:spectral}
\end{table}

The eigenvalue distributions, which control the \emph{shape} of the coupling landscape (how many strong coupling modes, how they decay, what the effective rank is), are nearly identical across all layers.
The raw matrices are completely unrelated: different orientations of the same shape.

This is not a property of GPT-2 alone.
It holds across three architecture families (GPT-2, SmolLM2, Qwen2.5), three scales (124M--774M), and independently trained models.
The coupling geometry is invariant; only its orientation varies with depth.

\begin{observation}[The rotation explanation]
Why are the spectra invariant while the matrices differ?
Each layer's response heads ($W_V$, $W_O$, FFN) update the residual stream, effectively \textbf{rotating the coordinate frame} for the next layer.
Layer $\ell$'s coupling $M_\ell$ is approximately layer~1's coupling \emph{rotated} by the cumulative effect of response heads $1$ through $\ell{-}1$.
Eigenvalues are preserved under rotation, a mathematical fact, which explains identical spectra with different eigenvectors.
The transformer stores $L$ rotated copies of one coupling rule.
\end{observation}

% ======================================================================
\section{The Orbital Response Network}

\subsection{Moving Coupling to the Residual Stream}

The spectral invariance motivates a simple architectural change: replace $2L$ independent coupling matrices with a single shared pair $A, B \in \R^{d \times d}$, applied at every layer:
\begin{equation}
    q^{(\ell)} = \text{Norm}(h^{(\ell)}) \cdot A, \qquad k^{(\ell)} = \text{Norm}(h^{(\ell)}) \cdot B
\end{equation}
where $h^{(\ell)}$ is the residual stream at layer $\ell$.
The coupling geometry $\M = AB^\top$ lives on the residual stream, part of the shared medium, not an independent per-layer component.

The per-layer response heads ($W_V$, $W_O$, FFN) remain independent.
They serve two roles:
\begin{enumerate}[nosep]
    \item \textbf{Content extraction:} determining what information flows through the coupling pattern.
    \item \textbf{Field rotation:} updating the residual stream so that the next application of the same $\M$ produces a different coupling pattern.
\end{enumerate}

This second role is critical.
The rotation is entirely implicit, learned by gradient descent, never imposed, but it is what makes the architecture rich.
One coupling rule applied in $L$ progressively rotated frames produces $L$ distinct coupling patterns.

\input{orn_architecture_diagram}

\subsection{What We Share vs What Prior Work Shares}

\begin{table}[h]
\centering\footnotesize
\begin{tabular}{@{}lcccc@{}}
\toprule
& \textbf{Q,K} & \textbf{V,O / FFN} & \textbf{Dynamic} & \textbf{Result} \\
\midrule
ALBERT~\cite{lan2020albert} & Shared & Shared & No & $-1.5\%$ \\
Univ.\ Transformer & Shared & Shared & No & Comp.\ gain \\
MoEUT~\cite{moeut2024} & Shared & Shared+MoE & No & Matches \\
Relaxed Rec.~\cite{mol2025} & +LoRA & +LoRA & No & Matches \\
MASA~\cite{masa2025} & Dict. & Dict. / --- & No & $-0.8\%$ \\
\textbf{ORN (ours)} & \textbf{Shared} & \textbf{Per-layer} & \textbf{Yes} & \textbf{Matches+} \\
\bottomrule
\end{tabular}
\caption{What is shared across layers. Only the ORN shares the empirically invariant component (coupling) while keeping response heads per-layer, enabling dynamic field-state evolution.}
\end{table}

\paragraph{MASA: the closest concurrent work.}
MASA~\cite{masa2025} is the most directly comparable approach, validated at real scale (110M--729M, plus Llama~3.1 8B post-hoc compression).
Each layer's projection matrix is reconstructed as a weighted sum of shared ``dictionary atoms'': $\hat{W}_\ell = \sum_s c_{\ell s} \cdot D_s$, where $D_s$ are shared atoms and $c_{\ell s}$ are per-layer scalar coefficients.
With $S = L/3$ atoms, MASA achieves 66.7\% attention parameter reduction at $< 1\%$ quality cost.

Three findings from MASA independently confirm our architectural choices:
\begin{enumerate}[nosep]
    \item \textbf{Q,K are more compressible than O.}
    MASA finds that ``compressing O introduces a bottleneck that harms language modelling.''
    Their best configuration (MASA-QKV) shares Q,K,V but keeps O per-layer, approaching the ORN's design of sharing coupling but not response.
    \item \textbf{The quality gap widens at scale.}
    MASA-QKVO loses 0.82\% at 729M, up from near-zero at 110M.
    Static scalar coefficients may not capture the depth-dependent variation that the ORN's field-state dynamics handle naturally.
    \item \textbf{No spectral analysis.}
    MASA notes ``more sophisticated clustering methods (e.g., spectral) are left for future work.''
    The spectral invariance that \emph{explains} why their sharing works is exactly our Experiment~0 finding.
\end{enumerate}

The fundamental difference is static vs dynamic: MASA reconstructs each layer's weights independently from a shared dictionary (no interaction between layers), while the ORN applies the same $\M$ to a field state that \emph{evolves}, each layer's coupling depends on what all previous layers computed.
MASA treats sharing as \textbf{compression} (represent $L$ matrices with fewer basis elements); the ORN treats it as \textbf{dynamics} (one generator iterated through an evolving field).

\subsection{The Semigroup Interpretation}

The shared $\M$ is not just a compressed weight matrix, it is a \textbf{generator} of a coupling semigroup.
The set of coupling patterns $\{\text{Attn}^{(1)}, \text{Attn}^{(2)}, \ldots, \text{Attn}^{(L)}\}$ is the image of the semigroup generated by $\M$ acting on the evolving residual field.
Each iteration produces a distinct coupling pattern because the field state has been transformed by the previous layer's response heads.

This replaces \emph{stored coupling} ($2Ld^2$ parameters) with \emph{orbital computation} ($2d^2$ parameters $\times$ $L$ iterations).
The semigroup's capacity, the richness of the coupling patterns it can generate, is determined by $\M$'s spectral structure and the response heads' ability to rotate the field between iterations.

\subsection{Perturbative Corrections}

The ORN includes a third component: context-gated perturbative corrections at each layer:
\begin{equation}
    h^{(\ell+1)} = h^{(\ell)}_{\text{orbital}} + g(h^{(\ell)}) \cdot \delta(h^{(\ell)})
\end{equation}
where $g$ is a learned gate (sigmoid, initialised near zero) and $\delta$ is a small correction MLP.
The gate opens as the model learns when context is informative enough to warrant a correction.
This adds instance-specific detail that the coupling geometry, which captures the \emph{average} relational structure, cannot provide.

% ======================================================================
\section{Results}

\subsection{V1: 91M Parameters on OpenWebText}

The V1 model is the first ORN trained at a scale comparable to GPT-2 Small (124M).

\begin{table}[h]
\centering\small
\begin{tabular}{@{}lcc@{}}
\toprule
& \textbf{ORN V1} & \textbf{GPT-2 Small} \\
\midrule
Total parameters & 91M & 124M \\
Coupling parameters & 524K (0.6\%) & ${\sim}15$M (12\%) \\
Response heads + FFN & 90.4M (99.4\%) & ${\sim}109$M (88\%) \\
Val loss (OWT) & 3.35 & 3.35 \\
Training data & 2.5B tokens OWT & 40B tokens WebText \\
Training cost & \$5.50 & N/A \\
\bottomrule
\end{tabular}
\caption{V1 ORN vs GPT-2 Small. The ORN matches GPT-2's validation loss with 27\% fewer total parameters and ${\sim}30\times$ fewer coupling parameters, trained on $16\times$ less data.}
\label{tab:v1}
\end{table}

\paragraph{Spectral crystallisation.}
$\M$'s effective rank drops from the initial random value to 19--20 and stabilises by 250M tokens, just 10\% of the total training budget.
The asymmetry ratio ($\|A\| / \|B\|$) decreases from 1.49 to 1.40 over training while individual singular values grow, indicating that $\M$ is sharpening its spectral structure (concentrating energy into fewer modes) rather than collapsing.
The 20 dominant eigenvectors of a $512 \times 512$ matrix, roughly 10K floating-point numbers, encode the full relational grammar the model uses.

\paragraph{Generated text.}
V1 produces grammatically correct, topically coherent multi-sentence text.
Outputs maintain subject-verb agreement, paragraph structure, and thematic consistency across sequences of 200+ tokens.

\subsection{Scaling: 1M to 50M Parameters}

\begin{table}[h]
\centering\small
\begin{tabular}{@{}rcccc@{}}
\toprule
\textbf{Scale} & \textbf{ORN} & \textbf{Transformer} & \textbf{$\Delta$} & \textbf{Coupling compression} \\
\midrule
1M & 5.125 & 5.231 & $-2.0\%$ & $L\times$ \\
5M & 3.564 & 3.453 & $+3.2\%$ & $L\times$ \\
10M & 2.887 & 2.805 & $+3.0\%$ & $12\times$ \\
50M & 2.117 & 2.128 & $-0.5\%$ & $12\times$ \\
\bottomrule
\end{tabular}
\caption{Val loss on TinyStories at matched total parameters. The ORN matches the transformer at every scale with $12\times$ fewer coupling parameters.}
\end{table}

\subsection{Frozen-M Transfer}

\paragraph{Cross-domain (50M, TinyStories-era).}
We trained a 50M ORN on OpenWebText, froze the coupling $\M$, and fine-tuned only the response heads on new domains.
For comparison, a transformer with frozen $W_Q$, $W_K$ was fine-tuned identically.

\begin{table}[h]
\centering\small
\begin{tabular}{@{}lcccc@{}}
\toprule
\textbf{Transfer} & \textbf{ORN zero-shot} & \textbf{ORN final} & \textbf{TF zero-shot} & \textbf{TF final} \\
\midrule
OWT $\to$ WikiText & 5.96 & \textbf{4.19} & 7.24 & 4.78 \\
OWT $\to$ Code (Python) & 3.89 & \textbf{1.94} & 5.93 & 2.39 \\
\bottomrule
\end{tabular}
\caption{Domain transfer with frozen coupling. The ORN's shared $\M$ transfers 1.3 nats better on WikiText and 2.0 nats better on Code at zero-shot. Advantages persist after fine-tuning.}
\end{table}

\paragraph{Cross-dataset (V1, 91M).}
The V1 frozen-$\M$ experiment provides the strongest evidence yet.
We took the crystallised $\M$ from V1 (trained on OpenWebText) and trained a fresh model on FineWeb-Edu with $\M$ completely frozen.
The frozen-$\M$ model tracks within 2\% of a fully trained model throughout training (val ${\sim}3.40$ vs 3.35).
The coupling geometry learned on web text, roughly 10K numbers (20 eigenvectors $\times$ 512 dimensions), encodes relational grammar that transfers without modification to a curated educational corpus.
This confirms the core hypothesis: coupling geometry is dataset-invariant.

\subsection{Multimodal Results}

We tested whether $\M$'s coupling geometry generalises across modalities, not just text domains.

\begin{table}[h]
\centering\small
\begin{tabular}{@{}lccc@{}}
\toprule
\textbf{Modality} & \textbf{ORN vs TF} & \textbf{M rank} & \textbf{Notes} \\
\midrule
Text (baseline) & $-0.5\%$ & 51 & Reference \\
Images (CIFAR-10) & $+0.4\%$ at 8K steps & 4 & Crossover: ORN overtakes TF \\
Music (MAESTRO) & $-3.9\%$ & 2 & Corrections compensate \\
Env.\ sounds (ESC-50) & $-3.1\%$ & 8 & Corrections compensate \\
\bottomrule
\end{tabular}
\caption{Multimodal results. ORN matches or beats the transformer across four modalities.}
\label{tab:multimodal}
\end{table}

\paragraph{Sub-linear rank scaling.}
Four modalities require only rank 4 in the shared $\M$, a 3.8$\times$ compression relative to the number of modalities.
Music coupling collapses to rank 2, consistent with the highly regular temporal structure of musical sequences.

\paragraph{Cross-modal transfer.}
A text-trained $\M$ works for images with only $+3.2\%$ loss, confirming that the coupling geometry captures relational structure general enough to span modalities.

\subsection{MoE Composition}

Mixture-of-experts composes orthogonally with the ORN (tested at 66M parameters on TinyStories):

\begin{center}\small
\begin{tabular}{@{}lcccc@{}}
\toprule
& \textbf{Params} & \textbf{Coupling} & \textbf{Val Loss} & \textbf{vs TF Dense} \\
\midrule
Transformer Dense & 22.4M & 1,573K & 2.586 & --- \\
ORN Dense & 20.9M & 131K & 2.576 & $-0.4\%$ \\
Transformer + MoE & 66.4M & 1,573K & 2.423 & $-6.3\%$ \\
ORN + MoE & 65.0M & 131K & 2.425 & $-6.2\%$ \\
\bottomrule
\end{tabular}
\end{center}

The MoE improvement (${\sim}6.3\%$) is identical for both attention types.
ORN + MoE matches Transformer + MoE with $12\times$ fewer coupling parameters.

\subsection{Inspectability}

Because $\M$ is a single shared object, it is directly inspectable from any checkpoint:
\begin{itemize}[nosep]
    \item Eigenspectrum reveals effective coupling rank (19--20 at V1 scale, 51 on TinyStories, 4 on images)
    \item Spectral crystallisation during training (rank sharpens and stabilises by 10\% of training)
    \item Q/K vector orbits show smooth rotation across layers (adjacent cosine ${\sim}0.9$)
    \item Coreference emerges in K-space (``Lily'' $=$ ``She'') without supervision
    \item Layer specialisation (local $\to$ structural $\to$ entity-focused) from one $\M$
\end{itemize}

None of these diagnostics are available from per-layer coupling matrices, where the coupling is distributed across $2L$ independent objects.

% ======================================================================
\section{M as a Koopman Operator}

The semigroup interpretation of $\M$ has a precise mathematical identity: $\M$ is a finite-dimensional \textbf{Koopman operator}~\cite{koopman1931hamiltonian}.

The Koopman framework, originating in Hamiltonian mechanics (1931), represents nonlinear dynamical systems through a linear operator acting on \emph{observables} (functions of the state) rather than on the state directly.
The key insight is that nonlinear dynamics in state space become linear dynamics in function space, at the cost of working in an infinite-dimensional space.
In practice, one truncates to a finite-dimensional approximation.

The ORN performs exactly this truncation.
The residual stream $h^{(\ell)}$ is the state; the query and key projections $h \cdot A$ and $h \cdot B$ are observables; $\M = AB^\top$ is the linear operator that maps between them.
The eigendecomposition of $\M$ then has direct dynamical meaning:

\begin{description}[nosep, leftmargin=1.5em]
    \item[Eigenvalues] are timescales.
    Large eigenvalues correspond to fast coupling modes (syntax, local agreement), small eigenvalues to slow modes (discourse structure, long-range coherence).
    \item[Eigenvectors] are independent dynamical modes.
    Each eigenvector defines a direction in representation space along which coupling evolves independently.
    V1's rank-20 $\M$ means the model uses 20 independent relational channels.
\end{description}

This identification connects the ORN to 90+ years of dynamical systems theory.
A transformer stores what happened (per-layer weight snapshots).
The ORN stores how things evolve (a dynamical operator whose spectrum encodes the grammar of relations).

The Koopman perspective also explains why $\M$ transfers across datasets: the operator captures the \emph{dynamics} of how tokens relate, not the \emph{content} of any particular dataset.
Relational grammar, the fact that subjects relate to verbs, that brackets match, that questions expect answers, is a dynamical invariant.

% ======================================================================
\section{CORN: Composable Orbital Response Network}

\subsection{The Prediction--Composability Tension}

The ORN is trained for next-token prediction: given context, predict the next token.
But prediction and \emph{compositional reasoning} are distinct capabilities.
A model can predict the next word in a sentence without being able to compose multi-step inferences.

The semigroup property, $T(s+t) = T(s) \circ T(t)$, the requirement that transitions compose, is precisely the mathematical statement of compositional reasoning: the result of a two-step process equals two one-step processes chained together.
Standard training does not enforce this.

\subsection{Time-Scaled Coupling}

CORN introduces a single scalar parameter $t$ that scales the coupling:
\begin{equation}
    q = h \cdot (t \cdot A), \qquad k = h \cdot (t \cdot B)
\end{equation}
No new parameters are added.
At $t = 1$, CORN reduces to the standard ORN.
At other values of $t$, the model explores different timescales of the coupling semigroup.

The semigroup consistency is enforced by a Chapman-Kolmogorov (CK) loss:
\begin{equation}
    \mathcal{L}_{\text{CK}} = \left\| T(s+t) - T(s) \circ T(t) \right\|
\end{equation}
which penalises violations of the composition property.

\subsection{Three-Phase Training}

CORN uses a three-phase training recipe:
\begin{description}[nosep, leftmargin=1.5em]
    \item[Phase 1:] Standard cross-entropy (CE) training. Establish prediction capability.
    \item[Phase 2:] CE + semigroup (SG) loss. Enforce compositional consistency while maintaining prediction.
    \item[Phase 3:] Freeze $\M$, fine-tune response heads only. Domain adaptation.
\end{description}

In initial experiments, the semigroup error drops 33\% with the CK loss, at a prediction cost of only 2.9\%.
The model learns to compose its coupling dynamics while barely sacrificing next-token prediction.

\subsection{The Complex Depth Hypothesis}

The time parameter $t$ need not be real.
We conjecture that real $t$ corresponds to prediction (advancing along the sequence), while imaginary $t$ corresponds to reasoning (exploring lateral connections without advancing).
This is a Wick rotation: the same mathematical structure that connects quantum mechanics to statistical mechanics connects prediction to reasoning.

At $t = i\tau$ (pure imaginary), the coupling becomes oscillatory rather than dissipative, enabling the model to revisit and recombine representations without committing to a next-token prediction.
This is speculative but mathematically well-defined, and connects to:
\begin{itemize}[nosep]
    \item Diffusion models~\cite{song2023consistency}, which reverse a dissipative process
    \item Deep-OSG~\cite{wu2023deep}, which parameterises operator semigroups for continuous-depth networks
    \item Consistency models~\cite{song2023consistency}, which enforce trajectory consistency
\end{itemize}

% ======================================================================
\section{How This Differs from Chain-of-Thought Thinking}

Current ``thinking'' models (o1, R1, QwQ) achieve reasoning by generating more tokens: the model writes out intermediate steps in natural language, each processed by the same forward pass.
This is effective but architecturally trivial: more tokens, same computation per token, no structural change.

CORN's approach is fundamentally different:

\begin{description}[nosep, leftmargin=1.5em]
    \item[Token-space reasoning (CoT):] Each reasoning step is a full forward pass through $L$ layers, producing a token that is appended to the context.
    Reasoning cost scales linearly with the number of reasoning steps, each step costs as much as any other token, and the intermediate steps must be expressible in natural language.

    \item[Latent-space reasoning (CORN):] Compositional reasoning occurs through semigroup dynamics in the residual stream.
    The time parameter $t$ controls the depth of composition.
    No tokens are generated; reasoning happens in the latent representation.
    The Koopman eigenmodes provide mathematical structure, each mode evolves independently at its own timescale, enabling principled decomposition of reasoning into parallel channels.
\end{description}

Recent work on latent reasoning supports this direction: COCONUT~\cite{coconut2024} replaces chain-of-thought tokens with continuous latent states, and Ouro/LoopLM~\cite{ouro2025} iterates through a looped transformer to refine representations without token generation.
Multi-token prediction~\cite{gloeckle2024multi} and DreamerV3~\cite{hafner2025dreamerv3} similarly move computation into latent space.

The key advantage is that CORN's latent reasoning has mathematical structure (Koopman modes, semigroup composition, spectral decomposition) rather than being an opaque latent state.
The eigenspectrum of $\M$ tells us \emph{how many} independent reasoning channels the model has, and which timescales they operate at.

% ======================================================================
\section{Discussion}

\subsection{The Implicit Semigroup}

The ORN's semigroup is emergent, not coded.
The implementation is simple parameter sharing: one $A$, one $B$, referenced by every layer.
We do not compute $\M^L$ or compose coupling matrices explicitly.
The semigroup arises because the per-layer response heads transform the field state between applications of $\M$, producing $L$ distinct patterns from one rule.

This is architecturally identical to ALBERT sharing attention weights, the only difference is \emph{what} we share.
ALBERT shares everything and loses ${\sim}1.5\%$; the ORN shares only the invariant component and loses nothing.

\subsection{Connection to the Path Integral}

The residual stream accumulates contributions from every layer:
$h_L = h_0 + \sum_\ell f_\ell(\M, h_\ell)$.
Each term $f_\ell$ involves a Boltzmann-weighted sum over all tokens (the softmax attention is algebraically identical to a partition function).
The double sum, over layers and over tokens, is a discrete path integral along an orbit in representation space.

The ORN does not compute the path integral in the compositional sense ($T^L$, summing over intermediate paths).
It uses the semigroup as a geometric scaffold for contextual refinement: the coupling geometry provides the grammar of which tokens relate; the response heads provide the content of what flows between them; the perturbative corrections add instance-specific detail.

A companion technical report (\textit{Path Integrals and the ORN}) explores this relationship in detail, including experiments that demonstrate exact physics computation (1D Ising thermodynamics at 3--4 significant figures) when the ORN is trained with explicit propagator supervision.

% ======================================================================

\section{Four Scales of Intelligence}

The ORN's architectural decomposition suggests a hierarchy of knowledge, each with its own training cost and update frequency:

\begin{description}[nosep, leftmargin=1.5em]
    \item[1. Coupling geometry ($\M$), the grammar of relations.]
    Rank 19--20 out of 512 dimensions at V1 scale.
    Crystallises by 250M tokens (10\% of training) on diverse text.
    Encodes ``subjects relate to verbs, brackets match, questions expect answers'', relational structure that holds across all domains and modalities.
    \textbf{Cost: billions of tokens. Freeze early. Train once.}

    \item[2. Response heads, domain skills.]
    Per-layer $W_V$, $W_O$, FFN.
    Determine what information flows through the coupling pattern and how the field state evolves.
    Domain-specific: a code response head extracts indentation structure; a legal response head extracts clause boundaries.
    \textbf{Cost: fine-tuning on domain data with frozen $\M$. Medium.}

    \item[3. Perturbative corrections, instance-specific detail.]
    Context-gated, per-layer.
    Sharpen the prediction from ``a noun goes here'' to ``specifically \emph{Lily}.''
    The gate controls how much correction each token needs based on accumulated context.
    \textbf{Cost: smallest component. Can be added or swapped cheaply.}

    \item[4. Tool calling, factual lookup.]
    Not stored in weights at all.
    ``The population of France is 68 million'' belongs in a retrieval system, not in matrix multiplications.
    A 125M ORN that is excellent at knowing \emph{when} to call a tool and \emph{how} to use the result outperforms a 7B model that has memorised facts badly.
    \textbf{Cost: zero training cost for new facts. Infinite updatability.}
\end{description}

\paragraph{Experimental evidence.}
Our transfer results directly test levels 1--2: frozen $\M$ (trained on web text) transfers to WikiText ($+1.3$ nats better than frozen per-layer Q,K), to Python code ($+2.04$ nats), and to FineWeb-Edu (within 2\% of full training).
$\M$ learned the relational grammar from web text and it applied to formal writing, programming, and curated educational content without retraining.

The perturbative ablation (0.449 nats gap) tests level 3: removing the corrections degrades prediction substantially, confirming they carry instance-specific detail that the geometry and response heads cannot provide alone.

% ======================================================================
\section{Limitations and Open Questions}

\begin{itemize}[nosep]
    \item \textbf{Scale.} V1 at 91M is encouraging but not definitive.
    The key question, whether the scaling advantage grows with model size, requires runs at 350M+ parameters.
    \item \textbf{CORN is early.} The semigroup loss reduces composition error by 33\%, but we have not yet demonstrated downstream reasoning improvements.
    The complex-$t$ hypothesis is mathematically grounded but experimentally untested.
    \item \textbf{Benchmarks.} ORN V1 (91M params, 0.6\% coupling) beats GPT-2 Small on HellaSwag (33.0\% vs 31.6\%) and scores 57.5\% on PIQA---the first downstream benchmark win for the architecture.
    Further benchmark evaluation across more tasks is ongoing.
    \item \textbf{Long context.} The ORN's coupling geometry is independent of sequence length, but we have not tested beyond 1024 tokens.
    \item \textbf{Multimodal scale.} Sub-linear rank scaling at toy scale is promising; whether it holds at production scale with real vision encoders is unknown.
\end{itemize}

% ======================================================================
\section{Theory Glossary}

\begin{description}[nosep, leftmargin=1.5em, style=nextline]

    \item[\textbf{Orbital Response Network (ORN)}]
    Architecture with three components: SharedM (coupling geometry), per-layer response heads, and perturbative corrections.
    Replaces $2L$ per-layer coupling matrices with a single shared coupling rule applied to an evolving residual stream.

    \item[\textbf{SharedM ($\M = AB^\top$)}]
    A single pair $A, B \in \R^{d \times d}$ shared across all layers.
    Defines the coupling geometry: which tokens attend to which.
    Asymmetric ($A \neq B$) is essential for autoregressive tasks.

    \item[\textbf{Response heads}]
    Per-layer $W_V$, $W_O$, and FFN.
    Extract content from the coupling pattern and rotate the residual stream for the next layer.
    Domain-specific; must be fine-tuned for new tasks.

    \item[\textbf{Perturbative corrections}]
    Context-gated per-layer refinements: $h \gets h_{\text{orbital}} + g(h) \cdot \delta(h)$.
    Gate initialised near zero; opens as training progresses.
    Carry instance-specific detail.

    \item[\textbf{Koopman operator}]
    A linear operator on observables that represents nonlinear dynamics.
    $\M$ is a finite-dimensional Koopman operator: eigenvalues are timescales, eigenvectors are independent dynamical modes.
    Origin: Koopman (1931)~\cite{koopman1931hamiltonian}.

    \item[\textbf{Semigroup property}]
    $T(s+t) = T(s) \circ T(t)$: transitions compose.
    The mathematical statement of compositional reasoning.
    The ORN's coupling patterns form an implicit semigroup through field-state evolution.

    \item[\textbf{Chapman-Kolmogorov (CK) loss}]
    Training loss enforcing semigroup consistency: $\mathcal{L}_{\text{CK}} = \|T(s+t) - T(s) \circ T(t)\|$.
    Penalises violations of the composition property.

    \item[\textbf{CORN (Composable ORN)}]
    Extension of ORN with time-scaled coupling ($t \cdot A$, $t \cdot B$) and CK loss.
    Trained for both prediction and composability.
    At $t=1$, reduces to standard ORN.

    \item[\textbf{Spectral crystallisation}]
    $\M$'s effective rank dropping and stabilising during training.
    V1: rank crystallises to 19--20 by 250M tokens (10\% of training).
    Indicates the model has discovered its relational grammar.

    \item[\textbf{Coupling compression}]
    Replacing $2Ld^2$ per-layer coupling parameters with $2d^2$ shared parameters.
    V1 achieves ${\sim}30\times$ compression (524K vs ${\sim}15$M).

    \item[\textbf{Fast modes / slow modes}]
    Large eigenvalues of $\M$ correspond to fast coupling modes (syntax, local agreement); small eigenvalues to slow modes (discourse structure, long-range coherence).
    The spectral gap between fast and slow modes determines the model's effective separation of timescales.

\end{description}

% ======================================================================

\begin{thebibliography}{99}

\bibitem{vaswani2017attention}
A.~Vaswani et al.
Attention Is All You Need.
\emph{NeurIPS}, 2017.

\bibitem{lan2020albert}
Z.~Lan et al.
ALBERT: A Lite BERT for Self-supervised Learning of Language Representations.
\emph{ICLR}, 2020.

\bibitem{dehghani2018universal}
M.~Dehghani et al.
Universal Transformers.
\emph{ICLR}, 2019.

\bibitem{moeut2024}
R.~Csord\'{a}s, K.~Irie, and J.~Schmidhuber.
MoEUT: Mixture-of-Experts Universal Transformers.
\emph{NeurIPS}, 2024.

\bibitem{mol2025}
H.~Bae et al.
Relaxed Recursive Transformers: Effective Parameter Sharing with Cross-Layer LoRA.
\emph{arXiv:2410.20672}, 2024.

\bibitem{bao2024qk}
F.~Bao, A.~Hataya, and R.~Karakida.
Self-attention Networks Localize When QK-eigenspectrum Concentrates.
\emph{ICML}, 2024.

\bibitem{aubry2024block}
T.~Aubry, C.~Meng, A.~Sugolov, and V.~Papyan.
Transformer Block Coupling and its Correlation with Generalization in LLMs.
\emph{arXiv:2407.07810}, 2024.

\bibitem{masa2025}
Y.~Wang et al.
MASA: Share Your Attention via Matrix-based Dictionary Learning.
\emph{arXiv:2508.04581}, 2025.

\bibitem{koopman1931hamiltonian}
B.~O.~Koopman.
Hamiltonian Systems and Transformation in Hilbert Space.
\emph{Proceedings of the National Academy of Sciences}, 17(5):315--318, 1931.

\bibitem{wu2023deep}
H.~Wu and S.~Chen.
Deep Operator Semigroups.
\emph{arXiv:2305.09780}, 2023.

\bibitem{song2023consistency}
Y.~Song, P.~Dhariwal, M.~Chen, and I.~Sutskever.
Consistency Models.
\emph{ICML}, 2023.

\bibitem{coconut2024}
S.~Hao et al.
Training Large Language Models to Reason in a Continuous Latent Space.
\emph{arXiv:2412.06769}, 2024.

\bibitem{ouro2025}
L.~Fan et al.
LoopLM: Folding LLMs into Loops.
\emph{arXiv}, 2025.

\bibitem{gloeckle2024multi}
F.~Gloeckle et al.
Better \& Faster Large Language Models via Multi-token Prediction.
\emph{ICML}, 2024.

\bibitem{hafner2025dreamerv3}
D.~Hafner et al.
Mastering Diverse Domains through World Models.
\emph{Nature}, 2025.

\end{thebibliography}

\end{document}
